% =====================================================================
% 05_partie_c.tex — requêtes SQL
% =====================================================================
\chapter{Partie C : requêtes SQL}
\label{chap:partieC}

\section{Objectifs de la partie C}

La partie C vérifie la capacité du schéma à répondre à des questions métier
variées : fiches clients, comptes, opérations, agrégats, collections,
signataires, échéances et suppressions conditionnelles. Elle s'appuie sur la
syntaxe SQL décrite dans la référence officielle \cite{oracleSQLRef}. Elle ne se limite pas à
des \code{SELECT *} : elle utilise le déréférencement, les jointures, les
sous-requêtes corrélées, les agrégats, les fonctions de classement et les
collections. Les requêtes sont numérotées de 6.1 à 6.25, certaines étant
décomposées en variantes \code{a}, \code{b} ou \code{c}.

\begin{objectifbox}[Ce que la partie C démontre]
La partie C prouve que le modèle objet reste interrogeable avec SQL. Les
\code{REF} se déréférencent, les \code{VARRAY} et les \emph{nested tables} se
transforment en lignes avec \code{TABLE(...)}, et les agrégats classiques
restent disponibles. Elle montre aussi que les résultats dépendent de l'état de
la base au moment de l'exécution.
\end{objectifbox}

\section{Les requêtes 6.1 à 6.25}

Le tableau~\ref{tab:requetes-C} reprend les thèmes, les techniques clés et les
résultats attendus au moment de l'exécution de la partie C, c'est-à-dire avant
les tests destructifs de la partie D. Il sert de carte de révision : chaque
ligne peut être développée en soutenance à partir du script
\code{partie\_C\_requetes.sql}.

{\small
\setlength{\tabcolsep}{3.5pt}
\begin{longtable}{@{}C{0.9cm}L{5.0cm}L{4.4cm}L{3.8cm}@{}}
\caption{Les requêtes de la partie C et leurs résultats attendus.}
\label{tab:requetes-C}\\
\toprule
\entete{N°} & \entete{Thème} & \entete{Technique clé} &
\entete{Résultat clé (au moment de C)} \\
\midrule
\endfirsthead
\toprule
\entete{N°} & \entete{Thème} & \entete{Technique clé} &
\entete{Résultat clé (au moment de C)} \\
\midrule
\endhead
\bottomrule
\endfoot
6.1 & Fiche client sans afficher les téléphones &
\code{SELECT} avec formats \code{COLUMN} & 6 clients. \\
6.2 & Comptes courants & \code{REF} déréférencé &
7 CC, dont CC1 à 2\,498\,500. \\
6.3a & Épargne triée par taux & \code{ORDER BY taux\_interet DESC} &
CE2 en tête (3,49). \\
6.3b & Comptes à terme & \code{SELECT} simple & 4 CT. \\
6.4 & CC du client 4 & Jointure implicite & CC4 et CC7. \\
6.5 & Tous les comptes du client 4 & \code{UNION} CC + CE + CT &
CC4, CC7, CE4. \\
6.6 & Clients + CC < 100\,000 & Jointure + filtre &
CC2, CC3, CC5, CC6. \\
6.7a & Clients avec au moins un CE & \code{EXISTS} / jointure &
Clients 1, 2, 3, 4. \\
6.7b & Clients avec un seul CE &
\code{GROUP BY ... HAVING COUNT(*) = 1} & Clients 2, 3, 4. \\
6.8 & Clients ayant opéré sur CC7 & Jointure \code{operation} &
Clients 2, 3, 4. \\
6.9a & Opérations CC7 par le propriétaire &
\code{o.ref\_client = cc.ref\_client} & Opérations du client 4. \\
6.9b & Opérations CC7 par non-propriétaires &
\code{o.ref\_client <> cc.ref\_client} & Opérations des clients 2 et 3. \\
6.10 & CC à solde négatif + titulaire + conseiller &
Jointures + \code{solde < 0} & CC2 : $-50\,000$. \\
6.11a & Nombre de CE par client & \code{GROUP BY} &
4 lignes : 2, 1, 1, 1 (clients 1 à 4). \\
6.11b & Idem avec clients sans CE & \code{LEFT JOIN} / \code{NVL} &
6 lignes, dont 0. \\
6.12 & Client avec le plus de CE & Sous-requête + \code{HAVING} &
Client 1 (2 CE). \\
6.13 & Par type : nombre et montant moyen &
\code{GROUP BY type\_op} & 20 opérations au moment de C. \\
6.14 & Total des soldes par agence & \code{UNION ALL} CC + CE &
A1 : 3\,478\,500 ; A2 : 700\,000 ; A3 : 158\,000. \\
6.15 & Top 3 clients par solde total &
\code{ORDER BY ... FETCH FIRST 3 ROWS} &
Client 1 : 3\,348\,500 ; client 4 : 830\,000 ; client 2 : 300\,000. \\
6.16a & Téléphones du client 1 & \code{TABLE(c.telephones)} &
3 éléments, 2 renseignés. \\
6.16b & Téléphones de tous les clients & \code{TABLE} + jointure &
6 clients. \\
6.17 & Numéros de téléphone de chaque client &
\code{TABLE(c.telephones)} + \code{NVL} &
6 clients, 3 éléments chacun, certains \code{<non renseigne>}. \\
6.18a & Nombre de signataires de CC7 &
\code{COUNT(*)} sur \code{TABLE(SELECT signataires ...)} & 3. \\
6.18b & Nombre de signataires par CC &
Sous-requête scalaire \code{COUNT(*)} sur \code{TABLE(cc.signataires)} &
CC1:1, CC2:1, CC3:2, CC4:3, CC5:0, CC6:0, CC7:3. \\
6.19a & Signataires de CC7 & \code{TABLE} & 3 lignes. \\
6.19b & Signataires de tous les CC & \code{TABLE} + jointure &
Tous les comptes. \\
6.19c & Idem avec curseur imbriqué & Curseur paramétré + \code{FOR} &
Affichage par compte. \\
6.20 & Signataires de CC7 + adresse &
\code{TABLE} + \code{s.ref\_client = REF(c)} &
Clients 2 (D), 3 (X), 5 (R). \\
6.21a & Clients par ville & \code{ORDER BY adresse.ville} &
Béjaïa, Alger, Oran. \\
6.21b & Clients de Béjaïa & \code{WHERE adresse.ville = 'Bejaia'} &
Clients 1 et 4. \\
6.22 & CT à échéance dans 30 jours &
\code{BETWEEN SYSDATE AND SYSDATE+30} & CT2 (5 jours), CT1 (20 jours). \\
6.23 & CC sans opération depuis 12 mois &
\code{NOT EXISTS} + sous-requête & CC5 et CC6. \\
6.24 & Opérations > 1\,000\,000 DA & Filtre \code{montant > 1000000} &
Opération de 1\,500\,000 sur CC1. \\
6.25 & Suppression CT expirés > 1 an, montant 0 &
\code{DELETE} + \code{ADD\_MONTHS} & CT4 supprimé. \\
\end{longtable}
}

\begin{attentionbox}[Chronologie des résultats]
La partie C s'exécute \emph{avant} les tests destructifs de la partie D. C'est
pourquoi 6.2 montre encore CC1, que 6.13 compte 20 opérations et que 6.24
trouve l'opération de 1\,500\,000~DA. Après la partie D, CC1 et ses cinq
opérations sont supprimés, et CE4 est supprimé par
\code{fermeCompte('CE4')}. Si les captures sont refaites après l'exécution
complète, ces résultats ne seront plus visibles. Il faut soit capturer la
sortie de C avant D, soit utiliser \code{\_travail/\_C.out} ou
\code{\_travail/\_FULL.out}.
\end{attentionbox}

\section{Déréférencement d'attributs}

Un attribut objet se déréférence avec la notation pointée. Il faut un alias de
table pour que Oracle sache à quelle ligne se rapporte l'attribut. Écrire
\code{ref\_parrain.num\_client} sans alias échoue avec \code{ORA-00904}. La
double notation \code{c.ref\_parrain.\allowbreak ref\_parrain.\allowbreak nom}
permet d'atteindre le parrain du parrain, ce qui sera exploité en partie E.

\begin{lstlisting}[caption={Déréférencement d'attributs objet.},label={lst:dereferencement}]
SELECT c.num_client, c.nom, c.adresse.ville, c.ref_conseiller.nom
  FROM client c;

SELECT c.num_client, c.nom,
       c.ref_parrain.nom             AS parrain,
       c.ref_parrain.ref_parrain.nom AS grand_parrain
  FROM client c;
\end{lstlisting}

\section{Interroger un \code{VARRAY}}

La fonction de table \code{TABLE(...)} transforme une collection en ensemble de
lignes. Pour un \code{VARRAY}, \code{COUNT(*)} compte les emplacements, tandis
que \code{COUNT(COLUMN\_VALUE)} ne compte que les valeurs non \code{NULL}.
\code{CARDINALITY} ne fonctionne pas sur un \code{VARRAY} dans cette version :
elle lève \code{ORA-00932}. Il faut donc utiliser \code{TABLE(...)} et
\code{COUNT(*)}.

\begin{lstlisting}[caption={Interrogation d'un \code{VARRAY} de téléphones.},label={lst:varray}]
SELECT c.num_client, t.COLUMN_VALUE AS telephone
  FROM client c, TABLE(c.telephones) t
 WHERE c.num_client = 1;

SELECT COUNT(*)              AS nb_elements,
       COUNT(t.COLUMN_VALUE) AS nb_renseignes
  FROM TABLE (SELECT telephones FROM client WHERE num_client = 1) t;
\end{lstlisting}

\section{Interroger une \emph{nested table}}

Une \emph{nested table} se manipule de la même manière avec \code{TABLE(...)}.
La jointure avec le client se fait par \code{s.ref\_client = REF(c)}. La
requête 6.20 combine ainsi la \emph{nested table} des signataires, le compte
courant et le client, pour afficher l'adresse complète des signataires de CC7.

\begin{lstlisting}[caption={Signataires de CC7 avec adresse.},label={lst:nested-table}]
SELECT c.num_client, c.nom,
       c.adresse.num_rue||' '||c.adresse.rue||', '||c.adresse.ville AS adresse,
       s.droit
  FROM cpt_courant cc, TABLE(cc.signataires) s, client c
 WHERE cc.num_compte = 'CC7'
   AND s.ref_client = REF(c)
 ORDER BY c.num_client;
\end{lstlisting}

\section{Agrégats, classements et \code{NOT EXISTS}}

Les agrégats classiques restent disponibles. La requête 6.14 additionne les
soldes des comptes courants et d'épargne par agence avec \code{UNION ALL}. La
requête 6.15 calcule le solde total d'un client en additionnant trois
sous-requêtes, puis classe les résultats et ne garde que les trois premiers
avec \code{FETCH FIRST 3 ROWS ONLY}.

\begin{lstlisting}[caption={Top 3 des clients par solde total.},label={lst:top3}]
SELECT c.num_client, c.nom,
       NVL((SELECT SUM(solde) FROM cpt_courant x WHERE x.ref_client = REF(c)), 0)
     + NVL((SELECT SUM(solde) FROM cpt_epargne x WHERE x.ref_client = REF(c)), 0)
     + NVL((SELECT SUM(solde) FROM cpt_terme   x WHERE x.ref_client = REF(c)), 0)
       AS solde_total
  FROM client c
 ORDER BY solde_total DESC
 FETCH FIRST 3 ROWS ONLY;
\end{lstlisting}

La requête 6.23 utilise \code{NOT EXISTS} plutôt que \code{NOT IN}. Ce choix
est plus robuste en présence de \code{NULL} : \code{NOT IN} peut produire un
résultat vide lorsque la sous-requête contient une valeur nulle. La
sous-requête corrélée utilise \code{REF(cc)} pour relier l'opération au compte
courant.

\begin{lstlisting}[caption={Comptes courants sans opération depuis douze mois.},label={lst:not-exists}]
SELECT cc.num_compte, cc.solde, cc.statut,
       (SELECT MAX(o.date_op) FROM operation o
         WHERE o.ref_compte = REF(cc)) AS derniere_operation
  FROM cpt_courant cc
 WHERE NOT EXISTS (SELECT 1 FROM operation o
                    WHERE o.ref_compte = REF(cc)
                      AND o.date_op > ADD_MONTHS(TRUNC(SYSDATE), -12))
 ORDER BY cc.num_compte;
\end{lstlisting}

\section{Filtres, échéances et suppression conditionnelle}

La requête 6.22 sélectionne les comptes à terme dont l'échéance tombe dans les
trente prochains jours. La requête 6.24 isole les opérations de plus d'un
million de dinars. La requête 6.25 est la seule requête destructive de la
partie C : elle supprime les comptes à terme expirés depuis plus d'un an et
dont le montant bloqué est nul. \code{CT4} est le cas de test prévu ; les
autres comptes à terme sont conservés.

\begin{lstlisting}[caption={Opérations de plus d'un million et suppression des comptes à terme expirés.},label={lst:filtres-suppression}]
SELECT o.id_op, o.date_op, o.montant, o.type_op,
       c.num_client, c.nom AS client, cc.num_compte
  FROM operation o, client c, cpt_courant cc
 WHERE o.ref_client = REF(c)
   AND o.ref_compte = REF(cc)
   AND o.montant > 1000000
 ORDER BY o.montant DESC;

DELETE FROM cpt_terme
 WHERE date_echeance < ADD_MONTHS(TRUNC(SYSDATE), -12)
   AND montant_bloque = 0;
COMMIT;
\end{lstlisting}

\section{Résultats à montrer en priorité}

Huit résultats suffisent à démontrer la maîtrise de la partie C. Ils couvrent
le déréférencement, les agrégats, les collections, les échéances et la
suppression conditionnelle.

\begin{enumerate}
  \item \textbf{6.10} : CC2 à $-50\,000$, avec titulaire et conseiller.
  \item \textbf{6.15} : top 3 clients (1 : 3\,348\,500 ; 4 : 830\,000 ;
        2 : 300\,000).
  \item \textbf{6.16a} : téléphones du client 1 (3 éléments, 2 renseignés).
  \item \textbf{6.18a} : CC7 a 3 signataires.
  \item \textbf{6.20} : signataires de CC7 avec adresse complète — clients 2
        (D), 3 (X), 5 (R).
  \item \textbf{6.22} : CT2 à 5 jours, CT1 à 20 jours.
  \item \textbf{6.24} : opération de 1\,500\,000~DA (avant suppression de CC1).
  \item \textbf{6.25} : CT4 supprimé, les autres CT restent.
\end{enumerate}

\begin{remarquebox}[Preuves disponibles]
Les sorties sont conservées dans \code{\_travail/\_C.out} et dans la section
« PARTIE\_C\_REQUETES.SQL » de \code{\_travail/\_FULL.out}. Les captures n°~13
à n°~24 du plan de l'annexe~\ref{ann:captures} couvrent les requêtes
prioritaires.
\end{remarquebox}

\section{Pièges de la partie C}

\begin{table}[htbp]
  \centering
  \small
  \caption{Pièges de la partie C et solutions.}
  \label{tab:pieges-partieC}
  \begin{tabularx}{\textwidth}{@{}L{5.0cm}L{4.8cm}X@{}}
    \toprule
    \entete{Piège} & \entete{Symptôme} & \entete{Solution} \\
    \midrule
    Déréférencement sans alias & \code{ORA-00904}. &
    Toujours écrire \code{c.adresse.ville}, \code{c.ref\_conseiller.nom}. \\
    \code{CARDINALITY} sur un \code{VARRAY} & \code{ORA-00932}. &
    \code{TABLE(...)} + \code{COUNT(*)}. \\
    \code{COUNT} sur un \code{VARRAY} en SQL & \code{ORA-00904}. & Idem. \\
    \code{NOT IN} avec \code{NULL} & Aucun résultat inattendu. &
    Préférer \code{NOT EXISTS}. \\
    Confondre \code{montant} signé et valeur absolue & Les totaux sont faux. &
    Les retraits sont négatifs ; utiliser \code{SUM(montant)} ou \code{ABS}. \\
    Capturer C après D & CC1, CE4 et l'opération de 1,5~M absents. &
    Capturer avant D ou utiliser \code{\_C.out}. \\
    \bottomrule
  \end{tabularx}
\end{table}

\section{Conclusion de la partie C}

La partie C démontre que le modèle relationnel-objet reste interrogeable avec
SQL standard. Les collections ne sont pas des boîtes noires : \code{TABLE(...)}
les transforme en lignes, et les jointures avec les \code{REF} permettent de
retrouver les objets référencés. Les agrégats et les fonctions de classement
restent utilisables. La seule précaution majeure est chronologique : les
résultats de C doivent être observés avant les tests destructifs de D.
